\subsection{PRESENTATIONS OF A GROUP}

Let G be a group and \(\mathbf{t} = (t_{i})_{i \in I}\) a family of elements of G. Let \(f_{t}\) be the unique homomorphism of the free group F(I) into G which maps i to \(t_{i}\) . The image of \(f_{t}\) is the subgroup generated by the elements \(t_{i}\) of G. The elements of the kernel of \(f_{t}\) are called the relators of the family t. t is called generating (resp. free, basic) if \(f_{t}\) is surjective (resp. injective, bijective).

Let G be a group. A presentation of G is an ordered pair \((\mathbf{t}, \mathbf{r})\) consisting of a generating family \(\mathbf{t} = (t_{i})_{i \in I}\) and a family \(\mathbf{r} = (r_{j})_{j \in J}\) of relators such that the kernel \(N_{t}\) of \(f_{t}\) is generated by the elements \(gr_{j}g^{-1}\) for \(g \in \mathrm{F}(\mathrm{I})\) and \(j \in J\) . It amounts to the same to say that \(N_{t}\) is the normal subgroup of F(I) generated by the \(r_{j}\) for \(j \in J\) (in other words, the smallest normal subgroup of F(I) containing the elements \(r_{j}\) ( \(j \in J\) ), cf.~§4, no. 4). By an abuse of language the generators \(t_{i}\) and the relations \(r_{j}(\mathbf{t}) = e\) are said to constitute a presentation of the group G.

Let I be a set and \(\mathbf{r} = (r_j)_{j \ensuremath{\in} J}\) a family of elements of the free group F(I). Let N(r) be the normal subgroup of F(I) generated by the \(r_j\) for \(j \ensuremath{\in} J\) . Let F(I, r) = F(I)/N(r) and \(\tau_i\) denote the class of \(i\) modulo N(r). The ordered pair \((\tau, \mathbf{r})\) with \(\tau = (\tau_i)_{i \ensuremath{\in} I}\) is a presentation of the group F(I, r); if G is a group and (t, r) is a presentation of G with t = (t\_i)\_\{i \ensuremath{\in} I\}, there exists a unique isomorphism \(u\) of F(I, r) onto G such that \(u(\tau_i) = t_i\) for all \(i \ensuremath{\in} I\) . The group F(I, r) is said to be defined by the generators \(\tau_i\) and the relators \(r_j\) , or by an abuse of language that it is defined by the generators \(\tau_i\) and the relations \(r_j(\tau) = e\) . When I = {[}1, n{]} and J = {[}1, m{]}, it is said that F(I, r) is defined by the presentation

\[
\langle \tau_ {1}, \dots , \tau_ {n}; r _ {1}, \dots , r _ {m} \rangle .
\]

If \(r_j = u_j^{-1}v_j\) with \(u_j\) and \(v_j\) in \(\mathrm{F(I)}\) , this presentation is equally denoted by the symbol

\[
\langle \tau_ {1}, \dots , \tau_ {n}; u _ {1} = v _ {1}, \dots , u _ {m} = v _ {m} \rangle .
\]

Examples. (1) The group defined by the presentation \(\langle \tau; \tau^q = e \rangle\) is cyclic of order \(q\) .

\begin{enumerate}
\def\labelenumi{(\arabic{enumi})}
\setcounter{enumi}{1}
\tightlist
\item
  The group defined by the presentation \(\langle x, y; xy = yx \rangle\) is isomorphic to \(\mathbf{Z} \times \mathbf{Z}\) .
\end{enumerate}

PROPOSITION 9. Let G be a group, \(\mathbf{t} = (t_i)_{i \in I}\) a generating family of G and \(\mathbf{r} = (r_j)_{j \in J}\) a family of relators of \(\mathbf{t}\) . The following conditions are equivalent:

\begin{enumerate}
\def\labelenumi{(\alph{enumi})}
\item
  The ordered pair \((\mathbf{t},\mathbf{r})\) is a presentation of G.
\item
  Let \(\mathbf{G}'\) be a group and \(\mathbf{t}' = (t_i')_{i \in \mathbf{I}}\) a family of elements of \(\mathbf{G}'\) . If \(r_j(\mathbf{t}') = e\) for all \(j \in \mathbf{J}\) , there exists a homomorphism \(u\) of \(\mathbf{G}\) into \(\mathbf{G}'\) such that \(u(t_i) = t_i'\) for all \(i \in \mathbf{I}\) .
\item
  Let \(\overline{\mathbf{G}}\) be a group and \(\overline{\mathbf{t}} = (\bar{t}_i)_{i\in \mathbf{I}}\) a generating family of \(\overline{\mathbf{G}}\) such that \(r_j(\overline{\mathbf{t}}) = e\) for all \(j\in \mathbf{J}\) . Every homomorphism of \(\overline{\mathbf{G}}\) into \(\mathbf{G}\) which maps \(\bar{t}_i\) to \(t_i\) for all \(i\in \mathbf{I}\) is an isomorphism.
\end{enumerate}

Let \(f\) denote the homomorphism of \(\mathbf{F}(\mathbf{I})\) into \(\mathbf{G}\) which maps \(i\) to \(t_i\) for all \(i \in \mathbf{I}\) and \(\mathbf{N}\) the kernel of \(f\) .

\begin{enumerate}
\def\labelenumi{(\alph{enumi})}
\item
  \(\Rightarrow\) (b): Suppose that \((\mathbf{t},\mathbf{r})\) is a presentation of G and let \(\mathbf{t}' = (t_i')_{i\in I}\) be a family of elements of a group \(\mathbf{G}'\) with \(r_j(\mathbf{t}') = e\) for all \(j\in J\) . Let \(f'\) be the homomorphism of F(I) into \(\mathbf{G}'\) characterized by \(f'(i) = t_i'\) for all \(i\in I\) . By hypothesis \(f'(r_j) = e\) for all \(j\in J\) and, as N is generated by the elements \(gr_jg^{-1}\) for \(j\in J\) and \(g\in \mathrm{F}(\mathrm{I}),f'(\mathrm{N}) = \{e\}\) . As the homomorphism \(f\colon \mathrm{F}(\mathrm{I})\to \mathrm{G}\) is surjective with kernel N, there exists a homomorphism \(u\colon \mathrm{G}\to \mathrm{G}'\) such that \(f' = u\circ f\) . Then \(u(t_i) = u(f(i)) = f'(i) = t_i'\) .
\item
  \(\Rightarrow\) (c): Suppose condition (b) holds. Let \(\mathbf{t} = (t_i)_{i\in I}\) be a generating family of a group G such that \(r_j(\mathbf{t}) = e\) for all \(j\in J\) and let \(v\) be a homomorphism of \(\overline{\mathbf{G}}\) into G such that \(v(t_i) = t_i\) for all \(i\in I\) . As the family \((t_i)_{i\in I}\) generates G, the homomorphism \(v\) is surjective. By property (b) there exists a homomorphism \(u\colon G\to \overline{\mathbf{G}}\) such that \(u(t_i) = \bar{t}_i\) for all \(i\in I\) . Then \(u(v(\bar{t}_i)) = \bar{t}_i\) for all \(i\in I\) and hence \(u \circ v\) is the identity on \(\overline{G}\) , which proves that \(v\) is injective. Hence \(v\) is an isomorphism and condition (c) holds.
\item
  \(\Rightarrow\) (a): Suppose condition (c) holds. Let \(t_i'\) be the canonical image of \(i\) in \(\mathbf{F}(\mathbf{I},\mathbf{r})\) and \(\mathbf{t}' = (t_i')_{i\in \mathbf{I}}\) ; then \(r_j(\mathbf{t}') = e\) for all \(j\in \mathbf{J}\) . As \(r_j(\mathbf{t}) = e\) for all \(j\in \mathbf{J}\) , there exists one and only one homomorphism \(v\) of \(\mathbf{F}(\mathbf{I},\mathbf{r})\) into \(\mathbf{G}\) such that \(v(\bar{t}_i) = t_i\) for all \(i\in \mathbf{I}\) . By (c), \(v\) is an isomorphism of \(\mathbf{F}(\mathbf{I},\mathbf{r})\) onto \(\mathbf{G}\) which transforms the presentation \((\mathbf{t}',\mathbf{r})\) of \(\mathbf{F}(\mathbf{I},\mathbf{r})\) into a presentation \((\mathbf{t},\mathbf{r})\) of \(\mathbf{G}\) .
\end{enumerate}

